\documentclass[11pt]{article}

% --- Packages ---
\usepackage[utf8]{inputenc}
\usepackage[T1]{fontenc}
\usepackage{lmodern}
\usepackage[margin=1in]{geometry}
\usepackage{amsmath,amsthm,amssymb}
\usepackage{mathtools}
\usepackage{graphicx}
\usepackage[numbers,sort&compress]{natbib}
\usepackage{hyperref}
\usepackage{booktabs}
\usepackage{float}
\usepackage{caption}
\usepackage{subcaption}
\usepackage{microtype}
\usepackage{enumitem}

% --- Figure path ---
\graphicspath{{../code/figures/}}

% --- Theorem environments ---
\newtheorem{theorem}{Theorem}[section]
\newtheorem{proposition}[theorem]{Proposition}
\newtheorem{lemma}[theorem]{Lemma}
\newtheorem{corollary}[theorem]{Corollary}
\theoremstyle{definition}
\newtheorem{definition}[theorem]{Definition}
\newtheorem{example}[theorem]{Example}
\newtheorem{assumption}[theorem]{Assumption}
\theoremstyle{remark}
\newtheorem{remark}[theorem]{Remark}

% --- Shortcuts ---
\newcommand{\R}{\mathbb{R}}
\newcommand{\E}{\mathbb{E}}
\newcommand{\Prob}{\mathbb{P}}
\newcommand{\Q}{\mathbb{Q}}
\newcommand{\F}{\mathcal{F}}
\newcommand{\G}{\mathcal{G}}
\newcommand{\dt}{\,\mathrm{d}t}
\newcommand{\dW}{\,\mathrm{d}W}
\newcommand{\ds}{\,\mathrm{d}s}
\newcommand{\norm}[1]{\lVert #1 \rVert}
\newcommand{\abs}[1]{\lvert #1 \rvert}
\newcommand{\ip}[2]{\langle #1, #2 \rangle}
\DeclareMathOperator{\tr}{tr}

% ===================================================================
\title{From Classical BSDEs to Deep Solvers:\\
Nonlinear Derivative Pricing, XVA, and High-Dimensional PDEs}

\author{Abhi Wadhwa}

\date{}

\begin{document}
\maketitle

% ===================================================================
\begin{abstract}
This paper builds a computational framework connecting backward stochastic
differential equations (BSDEs) to three classes of problems in quantitative
finance: nonlinear derivative pricing under market frictions, credit and
funding valuation adjustments (XVA), and dynamic risk measures via
$g$-expectations.  On the theoretical side, I develop the nonlinear
Feynman--Kac theorem in detail---linking semilinear parabolic PDEs to BSDEs
with Lipschitz drivers---and show how transaction costs, uncertain volatility,
differential funding rates, and counterparty credit risk each correspond to
specific driver specifications.  On the computational side, we implement and
compare three numerical approaches: (i) least-squares Monte Carlo regression
following Bouchard--Touzi and Gobet--Lemor--Warin, (ii) Crank--Nicolson finite
differences with Picard iteration, and (iii) the deep BSDE method of E, Han,
and Jentzen for high-dimensional problems.  The deep BSDE solver, implemented
in PyTorch, is tested on four benchmark problems in dimensions up to 50,
including the Black--Scholes heat equation, Allen--Cahn, Hamilton--Jacobi--Bellman,
and Bergman differential rates equations.  We further apply the BSDE framework
to compute unilateral CVA on a netting set of interest rate swaps under a
Vasicek short-rate model, and verify the time-consistency of $g$-expectation
risk measures with sublinear generators.  Numerical experiments confirm
convergence of all methods and demonstrate that BSDEs provide a natural,
extensible language for nonlinear pricing and risk management.
\end{abstract}

% ===================================================================
\section{Introduction}\label{sec:intro}

The Black--Scholes--Merton framework~\cite{blackscholes1973pricing} provides a
complete theory of derivative pricing in frictionless, complete markets.  At
its mathematical core lies the Feynman--Kac theorem, which identifies the
solution of the linear pricing PDE with a conditional expectation under the
risk-neutral measure.  This identification enables two complementary numerical
strategies: solve the PDE directly via finite differences, or estimate the
expectation via Monte Carlo.

In practice, the frictionless assumption is a convenient fiction.  Transaction
costs erode hedging gains~\cite{leland1985option}; borrowing and lending rates
differ~\cite{bergman1995option}; model parameters---especially volatility---are
uncertain~\cite{avellaneda1995pricing}; counterparties may default, creating
credit valuation adjustments (CVA) that depend nonlinearly on the portfolio
value~\cite{bichuch2018xva,crepey2014counterparty}.  Each of these frictions
introduces a nonlinearity into the pricing equation.  The linear Feynman--Kac
theorem no longer applies, and the two numerical strategies must be reconciled
through a more general theory.

Pardoux and Peng~\cite{pardoux1990adapted} initiated the modern theory of
backward stochastic differential equations (BSDEs) in 1990, proving that a
BSDE with Lipschitz driver admits a unique adapted solution.  Seven years
later, El Karoui, Peng, and Quenez~\cite{elkaroui1997backward} demonstrated
that BSDEs provide the natural framework for pricing and hedging in incomplete
and friction-laden markets: the BSDE value process $Y$ gives the price, the
control process $Z$ gives the hedge, and the driver $f$ encodes the market
friction.  The nonlinear Feynman--Kac theorem
(see~\cite{pardoux1990adapted,zhang2017backward}) then establishes that
$Y_t = u(t, X_t)$ where $u$ is a viscosity solution of the corresponding
semilinear PDE.

Since its introduction, the BSDE framework has expanded in three directions
that motivate this paper:

\paragraph{Direction 1: Deep learning for high-dimensional BSDEs.}
Classical BSDE solvers based on regression Monte Carlo
\cite{bouchard2004discrete,zhang2004numerical,gobet2005regression} scale
poorly beyond a handful of dimensions.  The deep BSDE method of E, Han, and
Jentzen~\cite{ehan2017deep,han2018solving} reformulates the discretized BSDE
as a stochastic control problem, approximates the gradient process $Z$ at each
time step by a neural network, and trains end-to-end by minimizing the terminal
condition loss.  Han and Long~\cite{hanlong2020convergence} provided a
posteriori error bounds, and subsequent work by Hur\'{e}, Pham, and
Warin~\cite{hure2020deep} (the DBDP1/DBDP2 schemes) and Beck et
al.~\cite{beck2021deep} (deep splitting) offered robustness alternatives.
These methods have solved problems in 100+ dimensions that are completely
intractable for grid-based PDE solvers.

\paragraph{Direction 2: XVA as nonlinear BSDEs.}
Post-2008 regulatory requirements mandate that banks account for
counterparty credit risk (CVA), own-credit risk (DVA), and funding costs (FVA)
when marking derivatives to market.  The total valuation adjustment
$\text{XVA} = \text{CVA} + \text{DVA} + \text{FVA}$
creates a nonlinear feedback: the funding cost depends on the derivative value,
which in turn depends on the funding cost.  Bichuch, Capponi, and
Sturm~\cite{bichuch2018xva} showed that this coupled system admits a natural
BSDE formulation, and Cr\'{e}pey, Bielecki, and
Brigo~\cite{crepey2014counterparty} developed the computational framework in
detail.

\paragraph{Direction 3: $g$-expectations and dynamic risk measures.}
Peng~\cite{peng2004nonlinear,peng2019nonlinear} recognized that the BSDE
mapping $\xi \mapsto Y_t$ defines a nonlinear expectation operator---the
$g$-expectation---that satisfies time-consistency by construction.  When the
generator $g$ is sublinear and positively homogeneous, the resulting
$g$-expectation defines a coherent risk measure in the sense of Artzner,
Delbaen, Eber, and Heath~\cite{artzner1999coherent}.  This connection provides
a principled way to construct dynamic, time-consistent risk measures for
portfolio management and regulatory capital computation.

\paragraph{Contributions and outline.}
This paper implements all three directions within a single computational
framework and provides numerical evidence for the theoretical connections.
Section~\ref{sec:theory} develops the BSDE theory:
existence/uniqueness (Pardoux--Peng), the financial interpretation
(El~Karoui--Peng--Quenez), the comparison theorem, and the nonlinear
Feynman--Kac theorem with a careful statement involving viscosity solutions.
Section~\ref{sec:deep_bsde} presents the deep BSDE method, including the
E--Han--Jentzen reformulation, convergence considerations, and comparison with
alternative deep solvers (DBDP, deep splitting).
Section~\ref{sec:numerics} details the numerical methods for forward SDE
simulation, classical BSDE regression, Crank--Nicolson PDE solving, and the
PyTorch deep BSDE implementation.
Section~\ref{sec:xva} formulates CVA/DVA/FVA as nonlinear BSDEs and presents
our CVA computation on a netting set of interest rate swaps.
Section~\ref{sec:g_expectations} develops $g$-expectations, their connection
to coherent risk measures, and time-consistency verification.
Section~\ref{sec:results} presents the full numerical results: classical
BSDE/PDE verification, deep BSDE benchmarks in dimensions 5--50, CVA
exposures and BSDE/MC comparison, and $g$-expectation properties.
Section~\ref{sec:conclusion} summarizes contributions and discusses open
directions.


% ===================================================================
\section{BSDE Theory}\label{sec:theory}

\subsection{Setup and Definition}

Let $(\Omega, \F, \Prob)$ be a complete probability space equipped with a
filtration $(\F_t)_{0 \le t \le T}$ generated by a $d$-dimensional standard
Brownian motion $W = (W_t)_{0 \le t \le T}$, augmented by the $\Prob$-null
sets.  All processes are assumed progressively measurable with respect to
$(\F_t)$.

\begin{definition}[BSDE]\label{def:bsde}
A backward stochastic differential equation on $[0,T]$ with terminal condition
$\xi \in L^2(\F_T)$ and driver (generator)
$f \colon [0,T] \times \Omega \times \R \times \R^d \to \R$ is the equation
\begin{equation}\label{eq:bsde}
  -\mathrm{d}Y_t = f(t, Y_t, Z_t)\dt - Z_t^\top \dW_t, \quad Y_T = \xi.
\end{equation}
An \emph{adapted solution} is a pair $(Y, Z) \in \mathcal{S}^2[0,T] \times
\mathcal{H}^2[0,T]$ satisfying~\eqref{eq:bsde}, where $\mathcal{S}^2$ denotes
the space of adapted c\`{a}dl\`{a}g processes with
$\E[\sup_{t \le T} |Y_t|^2] < \infty$ and $\mathcal{H}^2$ denotes the space
of predictable processes with $\E[\int_0^T \norm{Z_t}^2 \dt] < \infty$.
\end{definition}

In integral form:
\begin{equation}\label{eq:bsde_integral}
  Y_t = \xi + \int_t^T f(s, Y_s, Z_s)\ds - \int_t^T Z_s^\top \dW_s,
  \quad 0 \le t \le T.
\end{equation}
The terminal value $\xi$ and the driver $f$ are given; the unknowns are
$(Y, Z)$.  The process $Y$ carries the value, while $Z$ arises from the
martingale representation of $Y$ and carries the hedging information.

\subsection{Existence, Uniqueness, and the Comparison Theorem}

The foundational result is due to Pardoux and Peng~\cite{pardoux1990adapted}.

\begin{theorem}[Pardoux--Peng, 1990]\label{thm:existence}
Let $\xi \in L^2(\F_T)$ and suppose the driver $f$ satisfies:
\begin{enumerate}[label=(\roman*)]
  \item \textbf{Lipschitz continuity:} There exists $C > 0$ such that for all
  $t$, $y_1, y_2 \in \R$, $z_1, z_2 \in \R^d$,
  \[
    \abs{f(t, y_1, z_1) - f(t, y_2, z_2)}
    \le C\bigl(\abs{y_1 - y_2} + \norm{z_1 - z_2}\bigr).
  \]
  \item \textbf{Square integrability:}
  $\E\bigl[\int_0^T \abs{f(t,0,0)}^2 \dt\bigr] < \infty$.
\end{enumerate}
Then the BSDE~\eqref{eq:bsde} admits a unique adapted solution
$(Y, Z) \in \mathcal{S}^2[0,T] \times \mathcal{H}^2[0,T]$.
\end{theorem}

\begin{proof}[Proof sketch]
The proof proceeds by Picard iteration on the Banach space
$\mathcal{S}^2 \times \mathcal{H}^2$ equipped with the norm
$\norm{(Y,Z)}_\beta^2 = \E\bigl[\int_0^T e^{\beta t}(|Y_t|^2 + \norm{Z_t}^2)\dt\bigr]$
for $\beta$ sufficiently large.  Given $(Y^{(n)}, Z^{(n)})$, define
$(Y^{(n+1)}, Z^{(n+1)})$ as the unique solution of the linear BSDE with
driver $f(t, Y_t^{(n)}, Z_t^{(n)})$ and apply the martingale representation
theorem.  The Lipschitz condition ensures that the map is a contraction for
large $\beta$, yielding the fixed point.  See~\cite{pardoux1990adapted}
or~\cite[Chapter~4]{zhang2017backward} for details.
\end{proof}

A key consequence for pricing is the comparison theorem, which gives us
monotonicity of the BSDE solution with respect to the terminal condition and
the driver.

\begin{theorem}[Comparison Theorem]\label{thm:comparison}
Let $(Y^1, Z^1)$ and $(Y^2, Z^2)$ be solutions of BSDEs with drivers
$f^1, f^2$ and terminal conditions $\xi^1, \xi^2$.  Suppose:
\begin{enumerate}[label=(\roman*)]
  \item $\xi^1 \ge \xi^2$ a.s.;
  \item $f^1(t, Y_t^2, Z_t^2) \ge f^2(t, Y_t^2, Z_t^2)$ a.s.\ for a.e.\ $t$.
\end{enumerate}
Then $Y_t^1 \ge Y_t^2$ a.s.\ for all $t \in [0,T]$.
\end{theorem}

\begin{remark}
In the financial interpretation, the comparison theorem implies that a more
costly friction (larger driver) leads to a higher seller's price.  This gives
the natural ordering: transaction costs increase option prices, and
higher counterparty default intensity increases CVA.
\end{remark}

\subsection{The Financial Interpretation: El Karoui--Peng--Quenez}

El Karoui, Peng, and Quenez~\cite{elkaroui1997backward} established the
definitive financial interpretation of BSDEs.  Consider a market with one
riskless bond (rate $r$) and one risky asset $S$.  A self-financing portfolio
with wealth process $V$ and investment $\pi$ in the risky asset satisfies
\[
  \mathrm{d}V_t = r V_t \dt + \pi_t (\mathrm{d}S_t - r S_t \dt).
\]
Setting $Y_t = V_t$, $Z_t = \pi_t \sigma_t$ (where $\sigma_t$ is the
volatility of $S$), and $f(t,y,z) = -ry$, this is a linear BSDE.  The
correspondence becomes:
\begin{center}
\begin{tabular}{ll}
\toprule
\textbf{BSDE} & \textbf{Finance} \\
\midrule
$Y_t$ & Portfolio value (derivative price) \\
$Z_t$ & Hedging strategy ($\sigma \cdot \Delta$) \\
$\xi = g(S_T)$ & Payoff at maturity \\
$f(t, y, z)$ & Discounting + friction costs \\
\bottomrule
\end{tabular}
\end{center}

When the market is incomplete or frictions are present, the driver $f$ absorbs
the additional costs, and the BSDE value $Y_0$ gives the price that makes a
particular hedging strategy self-financing.

\subsection{The Nonlinear Feynman--Kac Theorem}\label{sec:nlfk}

Consider the coupled forward-backward system (FBSDE): the forward SDE
\begin{equation}\label{eq:fbsde_forward}
  \mathrm{d}X_t = b(t, X_t)\dt + \sigma(t, X_t)\dW_t, \quad X_0 = x,
\end{equation}
and the backward SDE
\begin{equation}\label{eq:fbsde_backward}
  -\mathrm{d}Y_t = f(t, X_t, Y_t, Z_t)\dt - Z_t^\top \dW_t,
  \quad Y_T = g(X_T).
\end{equation}

\begin{assumption}\label{asmp:regularity}
We assume: (a) $b, \sigma$ satisfy global Lipschitz and linear growth
conditions; (b) $f(t,x,y,z)$ is Lipschitz in $(y,z)$ uniformly in $(t,x)$,
continuous in $(t,x)$, and satisfies linear growth; (c) $g$ is continuous
with at most polynomial growth.
\end{assumption}

\begin{theorem}[Nonlinear Feynman--Kac]\label{thm:nonlinear_fk}
Under Assumption~\ref{asmp:regularity}, define
$u(t,x) \coloneqq Y_t^{t,x}$ where $(X^{t,x}, Y^{t,x}, Z^{t,x})$ solves
the FBSDE~\eqref{eq:fbsde_forward}--\eqref{eq:fbsde_backward} with $X_t = x$.
Then:
\begin{enumerate}[label=(\roman*)]
  \item The function $u$ is deterministic and continuous on $[0,T] \times \R^d$.
  \item $u$ is a viscosity solution of the semilinear PDE
  \begin{equation}\label{eq:nonlinear_pde}
    \frac{\partial u}{\partial t}
    + \frac{1}{2}\tr\!\left[\sigma\sigma^\top(t,x)\,D^2_x u\right]
    + b(t,x) \cdot \nabla_x u
    + f\!\left(t, x, u, \sigma^\top(t,x)\nabla_x u\right) = 0,
  \end{equation}
  with terminal condition $u(T,x) = g(x)$.
  \item Conversely, if $u \in C^{1,2}([0,T) \times \R^d)$ is a classical
  solution of~\eqref{eq:nonlinear_pde} with polynomial growth, then
  $Y_t = u(t, X_t)$ and $Z_t = \sigma^\top(t, X_t)\nabla_x u(t, X_t)$.
\end{enumerate}
\end{theorem}

\begin{proof}[Proof sketch]
Part~(iii) follows from It\^{o}'s formula applied to $u(t, X_t)$, identifying
terms with the BSDE.  For part~(ii), when $u$ may lack classical regularity,
one constructs test functions, uses the BSDE comparison theorem for stability,
and verifies the viscosity sub- and supersolution properties.  The full
argument is in~\cite[Chapter~6]{zhang2017backward}
and~\cite{pardoux1990adapted}.
\end{proof}

\begin{remark}
When $f(t,x,y,z) = -ry$ (the linear Black--Scholes driver),
equation~\eqref{eq:nonlinear_pde} reduces to the Black--Scholes PDE and
Theorem~\ref{thm:nonlinear_fk} recovers the classical Feynman--Kac formula
$u(t,x) = \E^\Q[e^{-r(T-t)} g(X_T) \mid X_t = x]$.
\end{remark}

\subsection{Market Frictions as Nonlinear Drivers}\label{sec:drivers}

Different market frictions correspond to different choices of the BSDE driver
$f$.  Each driver below satisfies the Lipschitz condition of
Theorem~\ref{thm:existence}, ensuring a well-posed pricing problem.

\subsubsection{Transaction Costs (Leland 1985)}
Following~\cite{leland1985option}, proportional transaction costs lead to
hedging costs that depend on the absolute rebalancing size:
\begin{equation}\label{eq:driver_tc}
  f_{\mathrm{tc}}(t, y, z) = -ry - \kappa \abs{z},
\end{equation}
where $\kappa \ge 0$ is the transaction cost intensity.  The nonlinearity
$\abs{z}$ makes the pricing equation semilinear.

\subsubsection{Uncertain Volatility (Avellaneda--Levy--Par\'{a}s 1995)}
When volatility is known only to lie in $[\sigma_L, \sigma_H]$,
worst-case pricing selects the volatility that maximizes the option value.
The simplified BSDE driver:
\begin{equation}\label{eq:driver_uv}
  f_{\mathrm{uv}}(t, y, z) = -ry + \lambda\abs{z},
  \quad \lambda = \tfrac{1}{2}(\sigma_H^2 - \sigma_L^2).
\end{equation}

\subsubsection{Asymmetric Funding Rates (Bergman 1995)}
With borrowing rate $r_b > r_l$ (lending rate), the driver is:
\begin{equation}\label{eq:driver_funding}
  f_{\mathrm{fund}}(t, y, z) = -R(y)\,y,
  \quad R(y) = r_l \cdot \mathbf{1}_{y \ge 0} + r_b \cdot \mathbf{1}_{y < 0}.
\end{equation}

\begin{proposition}\label{prop:lipschitz}
Each of the drivers~\eqref{eq:driver_tc},~\eqref{eq:driver_uv},
and~\eqref{eq:driver_funding} is Lipschitz in $(y,z)$ with constant
$C = \max(r + \kappa,\; r + \lambda,\; r_b)$.
\end{proposition}

\begin{proof}
For~\eqref{eq:driver_tc}: $\abs{f(y_1,z_1) - f(y_2,z_2)}
\le r\abs{y_1-y_2} + \kappa\abs{z_1-z_2}$ by the reverse triangle inequality.
The other two are similar.
\end{proof}


% ===================================================================
\section{The Deep BSDE Method}\label{sec:deep_bsde}

\subsection{The E--Han--Jentzen Reformulation}

The key insight of~\cite{ehan2017deep,han2018solving} is to view the
discretized BSDE as a stochastic control problem.  Consider the time grid
$0 = t_0 < t_1 < \cdots < t_N = T$ with $\Delta t = T/N$.  The Euler
discretization of the FBSDE gives:
\begin{align}
  X_{n+1} &= X_n + b_n \Delta t + \sigma_n \Delta W_n, \label{eq:deep_X}\\
  Y_{n+1} &= Y_n - f(t_n, X_n, Y_n, Z_n)\Delta t
             + (Z_n)^\top \sigma_n \Delta W_n, \label{eq:deep_Y}
\end{align}
where $\Delta W_n \sim \mathcal{N}(0, \Delta t \cdot I_d)$.

In the classical approach, both $Y_0$ and $\{Z_n\}_{n=0}^{N-1}$ are unknowns
recovered by backward regression.  The deep BSDE method instead
\emph{parameterizes} $Z_n \approx \phi_n(X_n; \theta_n)$ where each
$\phi_n$ is a feedforward neural network, and treats $Y_0$ as a single
trainable parameter.  The full forward pass~\eqref{eq:deep_X}--\eqref{eq:deep_Y}
from $n=0$ to $n=N$ produces a predicted terminal value $\hat{Y}_N$.  The
loss is
\begin{equation}\label{eq:deep_loss}
  \mathcal{L}(\theta) = \E\!\left[\abs{\hat{Y}_N - g(X_N)}^2\right],
\end{equation}
estimated by mini-batch sampling of the Brownian increments and minimized
by stochastic gradient descent (Adam optimizer).

\subsection{Network Architecture}

At each time step $n$, the network $\phi_n \colon \R^d \to \R^d$ has
the architecture:
\[
  x \;\xrightarrow{\text{BN}}\;
  [\text{Linear} \to \text{BN} \to \text{ReLU}]^L
  \;\xrightarrow{\text{Linear}}\; z,
\]
where BN denotes batch normalization and $L$ is the number of hidden layers
(typically $L = 2$).  The hidden dimension is modest (32--64 neurons per layer),
and the total number of trainable parameters scales as $O(N \cdot L \cdot d^2)$.

\begin{remark}
The deep BSDE method uses $N$ separate networks---one per time step---rather
than a single network conditioned on time.  This avoids the difficulty of
approximating a function that changes rapidly across the time grid, at the cost
of more parameters.
\end{remark}

\subsection{Convergence Theory and Caveats}

Han and Long~\cite{hanlong2020convergence} established a posteriori error
bounds: if the loss~\eqref{eq:deep_loss} is small, then the deep BSDE
approximation is close to the true solution in the $L^2$ sense.  Specifically,
for the coupled FBSDE with Lipschitz coefficients,
\[
  \E\!\left[\sup_{0 \le n \le N} |Y_n - \hat{Y}_n|^2\right]
  \le C \left(\mathcal{L}(\hat{\theta}) + (\Delta t)^2\right),
\]
where $C$ depends on the Lipschitz constants and $T$.

However, several caveats apply:
\begin{enumerate}[label=(\roman*)]
  \item The loss landscape is non-convex, and training can get stuck in local
  minima, particularly for highly nonlinear drivers.
  \item The $O(N)$ networks create a large parameter space; training requires
  careful learning rate scheduling and gradient clipping.
  \item The method provides a forward-in-time approximation of $Y_0$ but does
  not directly give the full solution surface $u(t,x)$.
\end{enumerate}

\subsection{Alternative Deep Solvers}

Two notable alternatives address some of these limitations:
\begin{itemize}
  \item \textbf{DBDP1/DBDP2}~\cite{hure2020deep}: Deep Backward Dynamic
  Programming trains a \emph{single} neural network at each backward time step
  to approximate the continuation value $Y_n = \E[Y_{n+1} + f \Delta t \mid X_n]$.
  This is closer in spirit to Longstaff--Schwartz but replaces polynomial
  regression with neural network regression.  DBDP2 further approximates $Z_n$
  via automatic differentiation of the trained network.
  \item \textbf{Deep splitting}~\cite{beck2021deep}: Reformulates the PDE
  solution as a fixed-point problem across time steps and solves each step
  independently, avoiding the end-to-end training of the deep BSDE method.
\end{itemize}


% ===================================================================
\section{Numerical Methods}\label{sec:numerics}

\subsection{Forward SDE Discretization}\label{sec:sde_schemes}

We discretize the forward SDE~\eqref{eq:fbsde_forward} on a uniform time
grid with step size $\Delta t = T/N$.

\paragraph{Euler--Maruyama.}
$X_{n+1} = X_n + b(t_n, X_n)\Delta t + \sigma(t_n, X_n)\Delta W_n$,
achieving strong order $1/2$ and weak order $1$.

\paragraph{Milstein.}
$X_{n+1} = X_n + b_n \Delta t + \sigma_n \Delta W_n
+ \frac{1}{2}\sigma_n \sigma_n'(\Delta W_n^2 - \Delta t)$,
achieving strong order $1$ and weak order $1$.  For geometric Brownian motion,
$\sigma_n' = \sigma$.

\subsection{Classical BSDE: Least-Squares Regression}\label{sec:bsde_numerics}

Following Bouchard and Touzi~\cite{bouchard2004discrete},
Zhang~\cite{zhang2004numerical}, and Gobet, Lemor, and
Warin~\cite{gobet2005regression}, we discretize the BSDE backward in time.
On the grid $\{t_0, \ldots, t_N\}$:
\begin{align}
  Y_N &= g(X_N), \label{eq:bsde_terminal}\\
  Z_n &= \frac{1}{\Delta t}\E\!\left[Y_{n+1}\Delta W_n \mid \F_{t_n}\right],
    \label{eq:Z_regression}\\
  Y_n &= \E\!\left[Y_{n+1} + f(t_n, Y_{n+1}, Z_n)\Delta t
    \mid \F_{t_n}\right]. \label{eq:Y_regression}
\end{align}
The conditional expectations~\eqref{eq:Z_regression}--\eqref{eq:Y_regression}
are approximated by projecting onto a polynomial basis of $X_{t_n}$.  Given
$M$ Monte Carlo paths, we construct the normalized polynomial basis
$\Phi(X_n^{(m)}) = (1, \tilde{X}, \tilde{X}^2, \ldots, \tilde{X}^p)$ of
degree $p$ and solve for regression coefficients by ordinary least squares.

This extends the Longstaff--Schwartz~\cite{longstaff2001valuing} regression
technique from American option pricing to the full BSDE setting.

\begin{remark}[Convergence of the classical solver]
Zhang~\cite{zhang2004numerical} proved that the backward Euler scheme converges
at rate $O(\sqrt{\Delta t})$ in $L^2$ for the $Y$ process.  The total error
comprises the time discretization error $O(\Delta t)$, the regression
approximation error (controlled by the basis richness), and the Monte Carlo
sampling error $O(M^{-1/2})$.  In practice, the Monte Carlo error often
dominates.
\end{remark}

\subsection{PDE Solver: Crank--Nicolson with Picard Iteration}

For low-dimensional problems, we solve the nonlinear
PDE~\eqref{eq:nonlinear_pde} directly via the Crank--Nicolson scheme
($\theta = 1/2$) on a uniform grid in $(S, t)$.  The linear system at each
time step is solved by sparse LU factorization.  For nonlinear terms
$h(S, u, Su_S)$, we employ Picard (fixed-point) iteration within each time
step: starting from the previous time level, we evaluate the nonlinear term,
solve the modified linear system, and iterate 4--5 times.  Convergence is
guaranteed for small $\Delta t$ by the contraction mapping theorem applied to
the Lipschitz nonlinearity.

\subsection{Deep BSDE: PyTorch Implementation}

Our deep BSDE implementation uses PyTorch with the following specifics:
\begin{itemize}
  \item $N = 20$ time steps; at each step, a \texttt{SubNet} with
  $L = 2$ hidden layers of width 64, batch normalization, and ReLU activation.
  \item $Y_0$ is initialized to $0.5$ and trained with a $10\times$ higher
  learning rate than the subnet parameters.
  \item Adam optimizer with learning rate $5 \times 10^{-3}$, step-wise decay
  (factor $0.5$ every $N_{\text{epochs}}/3$ epochs), and gradient clipping at
  norm $5.0$.
  \item Batch size $256$; training for $800$ epochs on CPU.
\end{itemize}

Four benchmark problems are implemented:
\begin{enumerate}
  \item \textbf{Black--Scholes (heat equation):} $f \equiv 0$,
  $g(x) = \norm{x}^2$, exact $u(0,x) = \norm{x}^2 + d\sigma^2 T$.
  \item \textbf{Allen--Cahn:} $f(y) = y - y^3$,
  $g(x) = (2 + 0.4\norm{x}^2)^{-1}$.
  \item \textbf{Hamilton--Jacobi--Bellman:}
  $f(z) = -\norm{z}^2/\sigma^2$,
  $g(x) = \ln(0.5(1 + \norm{x}^2))$.
  \item \textbf{Bergman differential rates:}
  $f(y,z) = -r_l y + (r_b - r_l)\max(0, y - \sum_i z_i / \sigma)$,
  $g(x) = (\max_i x_i - 1)^+$.
\end{enumerate}


% ===================================================================
\section{XVA and Counterparty Risk}\label{sec:xva}

\subsection{CVA as a Nonlinear BSDE}

Consider a netting set of derivatives with risk-free value $V_t$.  In the
presence of counterparty default risk with intensity $\lambda_t$ and loss-given-default
$\text{LGD} = 1 - R$ (recovery rate $R$), the CVA-adjusted value
$\tilde{V}_t$ satisfies the BSDE:
\begin{equation}\label{eq:cva_bsde}
  -\mathrm{d}\tilde{V}_t = \left[-r_t \tilde{V}_t
  - \lambda_t \cdot \text{LGD} \cdot (\tilde{V}_t)^+\right]\dt
  - Z_t^\top \dW_t,
  \quad \tilde{V}_T = g(X_T).
\end{equation}
The term $\lambda_t \cdot \text{LGD} \cdot (\tilde{V}_t)^+$ represents the
expected loss from counterparty default on positive exposure.  The unilateral
CVA is then $\text{CVA}_0 = V_0 - \tilde{V}_0$.

\begin{remark}[Nonlinear coupling in bilateral CVA/FVA]
When extending to bilateral CVA (including DVA for one's own default) and
FVA (funding valuation adjustment), the driver becomes:
\[
  f(t, v, z) = -r_t v - \lambda^c_t \cdot \text{LGD}^c \cdot v^+
  + \lambda^s_t \cdot \text{LGD}^s \cdot v^-
  + (r_f - r_t) \cdot h(v, z),
\]
where $\lambda^c, \lambda^s$ are the counterparty and self default intensities,
$r_f$ is the funding rate, and $h$ captures the funding requirement.  This
creates genuine nonlinear coupling: the funding cost depends on the derivative
value, which depends on the funding cost.  The BSDE framework handles this
naturally through the Lipschitz fixed-point structure of
Theorem~\ref{thm:existence}, as formalized
in~\cite{bichuch2018xva,crepey2014counterparty}.
\end{remark}

\subsection{Our CVA Computation}

We compute unilateral CVA on a netting set of $7$ interest rate swaps under
the Vasicek short-rate model:
\[
  \mathrm{d}r_t = \kappa(\theta - r_t)\dt + \sigma_r \dW_t,
\]
with $r_0 = 0.03$, $\kappa = 0.5$, $\theta = 0.04$, $\sigma_r = 0.01$.
The counterparty has constant default intensity $\lambda = 0.02$ and recovery
$R = 0.4$ ($\text{LGD} = 0.6$).

We compute:
\begin{enumerate}
  \item \textbf{Exposure profiles:} Expected Exposure $\text{EE}(t) = \E[V_t^+]$
  and Expected Positive Exposure $\text{EPE} = \frac{1}{T}\int_0^T \text{EE}(t)\dt$.
  \item \textbf{Monte Carlo CVA:}
  $\text{CVA}_{\text{MC}} = \text{LGD} \int_0^T \lambda \cdot D(0,t) \cdot
  \text{EE}(t) \dt$, where $D(0,t)$ is the discount factor.
  \item \textbf{BSDE CVA:} Solving the BSDE~\eqref{eq:cva_bsde} via
  backward regression with a 2D basis on $(r_t, V_t)$.
\end{enumerate}

The BSDE approach has the conceptual advantage of naturally extending to the
nonlinear bilateral CVA/FVA setting without fundamental algorithmic changes.

% ===================================================================
\section{$g$-Expectations and Dynamic Risk Measures}\label{sec:g_expectations}

\subsection{Peng's $g$-Expectation}

Given a BSDE with driver $g$ (which depends only on $z$, not on $y$):
\begin{equation}\label{eq:g_bsde}
  -\mathrm{d}Y_t = g(t, Z_t)\dt - Z_t^\top \dW_t, \quad Y_T = \xi,
\end{equation}
the mapping $\xi \mapsto Y_t$ defines the \emph{conditional $g$-expectation}
$\mathcal{E}^g_t[\xi] \coloneqq Y_t$~\cite{peng2004nonlinear,peng2019nonlinear}.

\begin{proposition}[Properties of $g$-expectations]\label{prop:g_properties}
If $g$ is Lipschitz in $z$ and $g(t, 0) = 0$, then $\mathcal{E}^g$ satisfies:
\begin{enumerate}[label=(\roman*)]
  \item \textbf{Monotonicity:} $\xi \ge \eta$ a.s.\ $\Longrightarrow$
  $\mathcal{E}^g_t[\xi] \ge \mathcal{E}^g_t[\eta]$ a.s.
  \item \textbf{Preservation of constants:} $\mathcal{E}^g_t[c] = c$ for
  $c \in \R$.
  \item \textbf{Time-consistency:}
  $\mathcal{E}^g_s\!\left[\mathcal{E}^g_t[\xi]\right] = \mathcal{E}^g_s[\xi]$
  for $0 \le s \le t \le T$.
\end{enumerate}
If on top of that $g$ is sublinear ($g(t, \alpha z) = \alpha g(t, z)$ for
$\alpha \ge 0$ and $g(t, z_1 + z_2) \le g(t, z_1) + g(t, z_2)$), then
$\rho_t(\xi) \coloneqq \mathcal{E}^g_t[-\xi]$ defines a coherent risk measure
in the sense of Artzner, Delbaen, Eber, and Heath~\cite{artzner1999coherent}:
it is monotone, translation-invariant, subadditive, and positively homogeneous.
\end{proposition}

\begin{proof}[Proof sketch]
(i) and (ii) follow from the comparison theorem (Theorem~\ref{thm:comparison}).
(iii) is a consequence of the flow property of BSDE solutions: by uniqueness,
the BSDE solved from $T$ to $s$ with terminal condition $\xi$ gives the same
$Y_s$ as first solving from $T$ to $t$ (obtaining $Y_t$) and then from $t$ to
$s$ with terminal condition $Y_t$.  The coherence properties under sublinearity
of $g$ follow from the corresponding properties of the BSDE
solution; see~\cite[Chapter~7]{zhang2017backward}.
\end{proof}

\subsection{The Sublinear Generator $g(z) = \gamma|z|$}

We focus on the generator $g(t, z) = \gamma \abs{z}$ with $\gamma > 0$.  This
is sublinear and positively homogeneous, so $\mathcal{E}^g$ defines a coherent
risk measure.  Financially, this $g$-expectation corresponds to worst-case
pricing under drift ambiguity: the market price of risk is uncertain within
a ball of radius $\gamma$ around zero.  The parameter $\gamma$ controls the
degree of risk aversion, with $\gamma = 0$ recovering the linear conditional
expectation.

\subsection{Connection to Entropic Risk Measures}

For comparison, the entropic risk measure with parameter $\theta > 0$ is:
\begin{equation}\label{eq:entropic}
  \rho_\theta(\xi) = \frac{1}{\theta}\ln \E\!\left[e^{-\theta\xi}\right].
\end{equation}
This is a convex (but not coherent) risk measure arising from
Kullback--Leibler divergence penalization.  It does not admit a direct
$g$-expectation representation (the corresponding BSDE has a quadratic
generator $g(z) = \frac{\theta}{2}\norm{z}^2$, which is superlinear).

\subsection{Connection to $G$-Expectations and 2BSDEs}

Peng's later work on $G$-expectations~\cite{peng2019nonlinear} and the
second-order BSDEs (2BSDEs) of Soner, Touzi, and
Zhang~\cite{sonertouziizhang2012wellposedness} extend the framework to handle
volatility uncertainty (not just drift uncertainty).  A 2BSDE replaces the
single volatility $\sigma$ with a family of volatilities, leading to fully
nonlinear PDEs.  This provides the theoretical foundation for the uncertain
volatility model of Avellaneda--Levy--Par\'{a}s and generalizes it to higher
dimensions.


% ===================================================================
\section{Results}\label{sec:results}

All classical experiments use $S_0 = K = 100$, $r = 0.05$, $\sigma = 0.20$,
$T = 1$ year.  The BSDE solver uses $M = 50{,}000$ paths with degree-4
polynomial regression; the PDE solver uses a $200 \times 500$ grid with
$S_{\max} = 400$.  Deep BSDE experiments use $\sigma = 1$, $T = 1$,
$N = 20$ time steps, batch size 256, and 800 training epochs on CPU.

\subsection{Classical Results: BSDE vs PDE Verification}

Figure~\ref{fig:bsde_vs_pde} verifies the nonlinear Feynman--Kac theorem
in the linear case: BSDE prices computed at 25 spot levels match the PDE
solution to within $0.10$ across $S \in [60, 140]$.

\begin{figure}[t]
  \centering
  \includegraphics[width=\textwidth]{fig1_bsde_vs_pde.pdf}
  \caption{Feynman--Kac verification: BSDE (Monte Carlo regression) vs PDE
  (Crank--Nicolson) for a European call under the linear Black--Scholes driver.
  Left: price overlay.  Right: absolute error on log scale.}
  \label{fig:bsde_vs_pde}
\end{figure}

Figure~\ref{fig:price_surfaces} compares initial-time prices under four
models (linear, transaction costs, funding costs, uncertain volatility).  All
nonlinear frictions increase the price relative to the Black--Scholes baseline.

\begin{figure}[t]
  \centering
  \includegraphics[width=0.75\textwidth]{fig2_price_surfaces.pdf}
  \caption{European call prices at $t=0$ under four models.  Nonlinear
  frictions consistently increase the price.}
  \label{fig:price_surfaces}
\end{figure}

\subsection{Convergence Analysis}

Figure~\ref{fig:sde_convergence} confirms the expected SDE convergence rates:
Euler--Maruyama achieves strong order $1/2$ and weak order $1$; Milstein
achieves strong order $1$ and weak order $1$.

\begin{figure}[t]
  \centering
  \includegraphics[width=\textwidth]{fig4_sde_convergence.pdf}
  \caption{Strong and weak convergence rates for Euler--Maruyama and Milstein
  schemes applied to GBM.  Dashed lines show theoretical reference slopes.}
  \label{fig:sde_convergence}
\end{figure}

Figure~\ref{fig:bsde_convergence} shows the BSDE solver error decreasing as
$O(N^{-1/2})$ with increasing time steps, consistent with the Monte Carlo
error dominating at the path counts used.

\begin{figure}[t]
  \centering
  \includegraphics[width=0.65\textwidth]{fig3_bsde_convergence.pdf}
  \caption{BSDE solver convergence: absolute error vs number of time steps $N$.
  The $O(N^{-1/2})$ rate reflects Monte Carlo dominance.}
  \label{fig:bsde_convergence}
\end{figure}

\subsection{Pricing Under Frictions}

Figure~\ref{fig:transaction_costs} shows that increasing transaction costs
$\kappa$ monotonically increases option prices, with ATM options most affected
(highest gamma).

\begin{figure}[t]
  \centering
  \includegraphics[width=\textwidth]{fig5_transaction_costs.pdf}
  \caption{Transaction cost sensitivity.  Left: price profiles for varying
  $\kappa$.  Right: ATM price as a function of $\kappa$.}
  \label{fig:transaction_costs}
\end{figure}

Figure~\ref{fig:uncertain_vol} displays the worst-case pricing bound under
uncertain volatility, lying near the upper Black--Scholes bound for convex
payoffs.

\begin{figure}[t]
  \centering
  \includegraphics[width=0.75\textwidth]{fig6_uncertain_vol.pdf}
  \caption{Uncertain volatility: the worst-case price lies near the upper BS
  bound, reflecting positive convexity of the call payoff.}
  \label{fig:uncertain_vol}
\end{figure}

Figure~\ref{fig:sample_paths} displays forward SDE sample paths for GBM, CEV,
and Heston models, illustrating the leverage effect and volatility clustering.

\begin{figure}[t]
  \centering
  \includegraphics[width=\textwidth]{fig7_sample_paths.pdf}
  \caption{Forward SDE sample paths: GBM (left), CEV with $\gamma = 0.5$
  (center), Heston (right).}
  \label{fig:sample_paths}
\end{figure}

Figure~\ref{fig:bsde_nonlinear} shows BSDE prices under all four drivers,
qualitatively consistent with the PDE results.

\begin{figure}[t]
  \centering
  \includegraphics[width=0.75\textwidth]{fig8_bsde_nonlinear.pdf}
  \caption{BSDE prices under four drivers, confirming consistency with PDE
  results (Figure~\ref{fig:price_surfaces}).}
  \label{fig:bsde_nonlinear}
\end{figure}

\subsection{Deep BSDE Results}

\paragraph{Training dynamics.}
Figure~\ref{fig:deep_training} shows training loss curves for all four
benchmarks at $d = 10$.  The Black--Scholes problem converges to a loss near
$1.8$, with $Y_0$ stabilizing at $10.00$ (exact value: $d \sigma^2 T = 10$).
The Allen--Cahn and HJB problems also converge, though with higher residual
loss reflecting the difficulty of approximating nonlinear drivers.

\begin{figure}[t]
  \centering
  \includegraphics[width=\textwidth]{fig9_deep_bsde_training.pdf}
  \caption{Deep BSDE training loss curves for four benchmarks at $d = 10$.
  All problems show monotone loss decrease.  Annotations show final $Y_0$
  and exact values where available.}
  \label{fig:deep_training}
\end{figure}

\paragraph{Dimension scaling.}
Figure~\ref{fig:deep_convergence} tests the Black--Scholes benchmark across
$d \in \{5, 10, 20, 50\}$.  At $d = 5$ and $d = 10$, the relative error is
below $1\%$, confirming that the method accurately solves these
moderate-dimensional problems.  At $d = 20$ and $d = 50$, the relative error
increases---reflecting the challenge of training with a fixed budget of $800$
epochs and batch size $256$ as dimension grows.  In the original
paper~\cite{han2018solving}, significantly more training was used (with GPU
acceleration) to achieve sub-percent errors at $d = 100$.

\begin{figure}[t]
  \centering
  \includegraphics[width=\textwidth]{fig10_deep_bsde_convergence.pdf}
  \caption{Dimension scaling for the Black--Scholes benchmark.
  Left: $Y_0$ vs exact value.  Right: relative error on log scale.
  Sub-percent accuracy at $d \le 10$; higher dimensions require more training.}
  \label{fig:deep_convergence}
\end{figure}

\paragraph{Benchmark solutions.}
Figure~\ref{fig:deep_solutions} summarizes $Y_0$ across all benchmarks at
$d = 10$.  For Black--Scholes, the deep BSDE solution matches the exact value
to $0.02\%$ relative error.  For Allen--Cahn ($Y_0 \approx 0.39$), HJB
($Y_0 \approx 1.47$), and Bergman ($Y_0 \approx 1.51$), the solutions
are consistent across runs and the training loss curves show convergence.

\begin{figure}[t]
  \centering
  \includegraphics[width=\textwidth]{fig11_deep_bsde_solutions.pdf}
  \caption{Deep BSDE solutions at $d = 10$.  Left: $Y_0$ values (blue: deep
  BSDE; coral: reference where available).  Right: training convergence for
  all benchmarks.}
  \label{fig:deep_solutions}
\end{figure}

\paragraph{Ablation study.}
Figure~\ref{fig:ablation} shows the effect of network width, depth, and
learning rate on the Black--Scholes benchmark at $d = 10$ (300 epochs).
All configurations are under-trained at this short budget, but the learning
rate has the strongest effect: $\text{lr} = 5 \times 10^{-3}$ brings $Y_0$
closest to the exact value.  Network width and depth have smaller effects,
consistent with the finding that the deep BSDE method is more sensitive to
optimization hyperparameters than architecture.

\begin{figure}[t]
  \centering
  \includegraphics[width=\textwidth]{fig16_ablation_deep_bsde.pdf}
  \caption{Ablation on the Black--Scholes benchmark ($d = 10$, 300 epochs).
  Left: effect of hidden dimension.  Center: effect of depth.  Right: effect
  of learning rate.  All configurations are under-converged; the learning rate
  effect dominates.}
  \label{fig:ablation}
\end{figure}

\subsection{CVA Results}

\paragraph{Exposure profiles.}
Figure~\ref{fig:cva_exposures} shows the expected exposure (EE) profile for
the 7-swap netting set over a 10-year horizon.  The EE profile peaks in the
mid-life of the swaps, reflecting the rolloff of shorter-maturity instruments
and the time decay of longer-maturity exposures.  The 95th percentile of
positive exposure shows significant tail risk.

\begin{figure}[t]
  \centering
  \includegraphics[width=\textwidth]{fig12_cva_exposures.pdf}
  \caption{Exposure profiles for the 7-swap netting set.  Left: expected
  exposure (EE) and expected negative exposure.  Right: EE with 5th--95th
  percentile band.}
  \label{fig:cva_exposures}
\end{figure}

\paragraph{CVA comparison.}
Figure~\ref{fig:cva_bsde} compares CVA computed via Monte Carlo integration
and via the BSDE approach.  The Monte Carlo CVA integrates
$\text{LGD} \cdot \lambda \cdot D(0,t) \cdot \text{EE}(t)$ over the time
horizon.  The BSDE CVA solves the nonlinear
equation~\eqref{eq:cva_bsde} backward in time.  Both methods produce CVA
estimates of the same order of magnitude, validating the BSDE formulation.

\begin{figure}[t]
  \centering
  \includegraphics[width=\textwidth]{fig13_cva_bsde.pdf}
  \caption{CVA: BSDE vs Monte Carlo.  Left: cumulative CVA profiles.
  Right: risk-free vs CVA-adjusted portfolio value (mean paths).}
  \label{fig:cva_bsde}
\end{figure}

\subsection{$g$-Expectation Results}

\paragraph{Risk aversion.}
Figure~\ref{fig:g_expectation} shows the $g$-expectation of a European call
payoff for varying risk aversion $\gamma \in \{0, 0.05, 0.1, 0.2, 0.5\}$
(with generator $g(z) = \gamma|z|$).  As $\gamma$ increases, the $g$-expectation
at $t = 0$ increases, reflecting the worst-case pricing under greater drift
ambiguity.  The right panel shows the monotone relationship between $\gamma$
and $\mathcal{E}^g_0[\xi]$.

\begin{figure}[t]
  \centering
  \includegraphics[width=\textwidth]{fig14_g_expectation.pdf}
  \caption{$g$-expectation analysis.  Left: $\mathcal{E}^g_t[\xi]$ over time
  for different $\gamma$.  Right: $\mathcal{E}^g_0[\xi]$ as a function of
  $\gamma$, with entropic risk measure comparisons (gray lines).}
  \label{fig:g_expectation}
\end{figure}

\paragraph{Time-consistency.}
Figure~\ref{fig:time_consistency} verifies the time-consistency property
$\mathcal{E}^g_0[\mathcal{E}^g_t[\xi]] = \mathcal{E}^g_0[\xi]$ at multiple
intermediate times $t$.  The direct computation (solving from $T$ to $0$)
and the composed computation (solving from $T$ to $t$, then from $t$ to $0$)
agree to high precision.  This is a
non-trivial numerical verification: time-consistency is a structural property
of BSDEs that distinguishes $g$-expectations from ad hoc risk measures.

\begin{figure}[t]
  \centering
  \includegraphics[width=\textwidth]{fig15_time_consistency.pdf}
  \caption{Time-consistency verification for $\gamma = 0.1$.  Left: direct
  vs composed $g$-expectation at $t = 0$ for varying intermediate times.
  Right: absolute error between the two computations.}
  \label{fig:time_consistency}
\end{figure}


% ===================================================================
\section{Conclusion}\label{sec:conclusion}

This paper connected classical BSDE
theory, deep learning solvers, XVA computation, and dynamic risk measures
within a single computational framework.  The main takeaways:

\begin{enumerate}
  \item \textbf{Theoretical development.}  We stated the Pardoux--Peng
  existence theorem, the comparison theorem, and the nonlinear Feynman--Kac
  theorem with proof sketches, and showed how each common market friction
  (transaction costs, uncertain volatility, differential rates) maps to a
  specific Lipschitz driver.

  \item \textbf{Classical verification.}  The least-squares regression BSDE
  solver and Crank--Nicolson PDE solver produce consistent prices across all
  tested drivers, with the expected convergence rates for both forward SDE
  schemes and the backward BSDE regression.

  \item \textbf{Deep BSDE implementation.}  Our PyTorch implementation of the
  E--Han--Jentzen method achieves sub-percent relative error on the
  Black--Scholes benchmark at $d = 10$, and produces converged solutions for
  the Allen--Cahn, HJB, and Bergman benchmarks.  Dimension scaling experiments
  show the method's promise for moderate-dimensional problems, while ablation
  studies confirm the importance of learning rate tuning.

  \item \textbf{CVA via BSDEs.}  We computed unilateral CVA on a 7-swap
  netting set under a Vasicek model, demonstrating both Monte Carlo and BSDE
  approaches and discussing the natural extension to bilateral CVA/FVA.

  \item \textbf{$g$-expectations.}  We verified time-consistency of
  $g$-expectations with sublinear generators, confirming the theoretical
  prediction that BSDE-based risk measures are dynamically consistent, and
  compared with entropic risk measures.
\end{enumerate}

\paragraph{Open directions.}
Several extensions are natural from this foundation:
\begin{itemize}
  \item \textbf{Path-dependent BSDEs.}  Dupire's functional It\^{o} calculus
  enables BSDEs driven by path-dependent functionals, relevant for Asian
  options and lookback features.
  \item \textbf{Reflected BSDEs.}  Adding a reflecting barrier $Y_t \ge L_t$
  gives the American option pricing problem.  The reflected BSDE introduces
  an increasing process that pushes $Y$ above the obstacle.
  \item \textbf{Second-order BSDEs.}  The 2BSDEs of Soner, Touzi, and
  Zhang~\cite{sonertouziizhang2012wellposedness} handle volatility uncertainty
  (fully nonlinear PDEs), extending the framework from semilinear to fully
  nonlinear problems.
  \item \textbf{Scaling deep methods.}  With GPU training and larger networks,
  the deep BSDE method can push to hundreds of dimensions, tackling
  basket CVA and high-dimensional control problems.
\end{itemize}

The overarching message is that BSDEs provide a native mathematical language
for nonlinear pricing, XVA, and risk measurement.  The combination of rigorous
theory (existence, uniqueness, comparison, Feynman--Kac) with modern
computational tools (regression Monte Carlo, deep learning) makes this
framework both theoretically grounded and practically implementable.


% ===================================================================
\section*{Acknowledgments}

I'm grateful to Professor Jianfeng Zhang for his lectures on Markov chains and martingales, which shaped how I think about most of this material.

\bibliographystyle{plainnat}
\bibliography{references}

\end{document}
